%The word �Abstract� should be centered 2? below the top of the page. 
%Skip one line, then center your name followed by the title of the 
%thesis/dissertation. Use as many lines as necessary. Centered below the 
%title include the phrase, in parentheses, �(Under the direction of  
%_________)� and include the name(s) of the dissertation advisor(s).
%Skip one line and begin the content of the abstract. It should be 
%double-spaced and conform to margin guidelines. An abstract should not 
%exceed 150 words for a thesis and 350 words for a dissertation. The 
%latter is a requirement of both the Graduate School and UMI's 
%Dissertation Abstracts International.
%Because your dissertation abstract will be published, please prepare and 
%proofread it carefully. Print all symbols and foreign words clearly and 
%accurately to avoid errors or delays. Make sure that the title given at 
%the top of the abstract has the same wording as the title shown on your 
%title page. Avoid mathematical formulas, diagrams, and other 
%illustrative materials, and only offer the briefest possible description 
%of your thesis/dissertation and a concise summary of its conclusions. Do 
%not include lengthy explanations and opinions.
%The abstract should bear the lower case Roman number ii (if you did not 
%include a copyright page) or iii (if you include a copyright page).

\begin{center}
\vspace*{52pt}
{\normalfont\textbf{ABSTRACT}}
\vspace{11pt}

\begin{singlespace}
ALEC L. PLOTKIN: Optimal Transport--Informed Mapping of CD8 T Cell Fate Across Tissues \\
(Under the direction of Natalie Stanley and Justin Milner)
\end{singlespace}
\end{center}

Single cell sequencing has revealed the profound heterogeneity of cell states, yet its destructive nature limits measurements of each cell to a single timepoint. Optimal transport (OT) methods have emerged as promising tools for inferring cell trajectories across multiple timepoints via probabilistic couplings between time-varying cell distributions. However, analyzing OT trajectories from single-cell data requires new frameworks. 
In this thesis, we introduce trajectory kernel mean embedding (TKME), which featurizes OT trajectories into compact summaries of cell states over time for downstream tasks like clustering. We apply OT to CD8 T cell responses during acute infection, which is a challenging system for other trajectory inference approaches. We demonstrate that OT recapitulates CD8 T cell population dynamics and that TKME partitions effector and memory cell lineages smoothly over time. 
We extend OT to model CD8 T cell migration between distinct biological compartments, and predicting that tissue-resident memory T cell fates are tightly coupled to their times of arrival. Finally, TKME enables compartment-wise comparison of circulating vs. resident memory differentiation, revealing context dependent roles for transcription factors TCF1 and AP-4, with T-bet as a universal driver of terminal effector differentiation. 
This work advances OT--informed modeling of cell trajectories through novel featurization and multi-compartment modeling, enabling biological discovery in complex experimental systems. 

\clearpage
